\chapter{动力学建模与观测机制}\label{ch:mechanism}

\section{微观生灭分支过程与生物学超传播机制}\label{sec:branching}

\subsection{经典 Galton--Watson 分支过程与代际传递}
在传染病暴发初期，当易感人群比例高度充沛（即易感者比例近乎恒定 $S/N \approx 1$）且尚未形成大规模群体免疫干预时，病原体的代际传播在微观尺度上可严格抽象为一个离散时间、连续状态空间的 Galton--Watson 随机分支过程 \citep{harris1963theory,jagers1975branching}。

设 $t \in \{0, 1, 2, \dots\}$ 表示感染代际离散时间索引。在第 $t-1$ 代，假设全社会存在 $I_{t-1} = n$ 名处于传染期的真实感染个体，构成第 $t$ 代传播的前瞻风险集（Risk Set）。假定这 $n$ 名个体在当前外部环境与生物学特性约束下相互独立地向外传播，其中第 $j$ 名感染者（$j = 1, 2, \dots, n$）在第 $t$ 代所引发的直接二代被感染人数记为离散随机变量 $X_{t,j}$。则第 $t$ 代全社会实际产生的新发真实感染总人数 $I_t$ 可表示为独立同分布随机变量之和：
\begin{equation}
I_t = \sum_{j=1}^n X_{t,j}. \label{eq:branching_sum}
\end{equation}

在经典均匀混合假定下，若假定每次接触传播服从均质泊松过程，则个体二代感染数 $X_{t,j} \sim \Poisson(R_t)$，其均值与方差均严格等于瞬时再生数 $R_t$。然而，大量的现代分子流行病学与流调溯源研究表明，均质泊松假定与真实呼吸道病毒传播的重尾特征存在显著偏差。

\subsection{个体传染力的 Gamma 连续混合模型：从微观连续积分到负二项 PMF}
为了刻画个体在生理排毒量、接触网络度分布以及社会行为上的极端异质性，本文遵循 Lloyd-Smith 等人 \citep{lloydsmith2005superspreading} 的经典生物数学设定，引入双层分层复合模型：
假定在特定时间代际 $t$ 内，全社会平均传播强度为 $R_t$；但对特定的个体 $j$ 而言，其实际具备的瞬时传染力为一个连续随机潜变量 $\nu_{t,j} \in (0, \infty)$。给定该个体的传染力 $\nu_{t,j}$，其产生的实际二代感染数服从条件泊松分布：
\begin{equation}
X_{t,j} \mid \nu_{t,j} \sim \Poisson(\nu_{t,j}), \quad \Pr(X_{t,j} = x \mid \nu_{t,j}) = \frac{\nu_{t,j}^x e^{-\nu_{t,j}}}{x!}, \quad x \in \{0, 1, 2, \dots\}.
\end{equation}
而在人群总体中，个体传染力 $\nu_{t,j}$ 的异质性分布由以形状参数 $k > 0$、尺度参数 $\theta = R_t / k$ 参数化的欧拉 Gamma 分布所刻画：
\begin{equation}
\nu_{t,j} \sim \GammaDist\left(k, \frac{R_t}{k}\right).
\end{equation}
其概率密度函数（PDF）严格由下式给出：
\begin{equation}
f_\nu(\nu; k, R_t) = \frac{1}{\theta^k \Gamma(k)} \nu^{k-1} e^{-\nu / \theta} = \frac{(k / R_t)^k}{\Gamma(k)} \nu^{k-1} \exp\left( -\frac{k \nu}{R_t} \right), \quad \nu > 0,
\end{equation}
其中 $\Gamma(k) \triangleq \int_0^\infty u^{k-1} e^{-u} du$ 为标准欧拉 Gamma 函数。

为了求解观测二代病例数 $X_{t,j}$ 的无条件概率质量函数（PMF），我们需要对未观测的连续潜变量 $\nu$ 在正半轴 $(0, \infty)$ 上进行边缘化全积分：
\begin{equation}
\Pr(X_{t,j} = x) = \int_0^\infty \Pr(X_{t,j} = x \mid \nu) \cdot f_\nu(\nu; k, R_t) \, d\nu.
\end{equation}

下面给出保姆级逐行代数积分演算步骤：

将条件泊松 PMF 与 Gamma 密度代入被积表达式：
\begin{equation}
\Pr(X_{t,j} = x) = \int_0^\infty \left[ \frac{\nu^x e^{-\nu}}{x!} \right] \cdot \left[ \frac{(k / R_t)^k}{\Gamma(k)} \nu^{k-1} \exp\left( -\frac{k \nu}{R_t} \right) \right] d\nu.
\end{equation}
提取所有与积分变量 $\nu$ 无关的常数因子，合并被积函数中的幂次项与指数幂次：
\begin{equation}
\Pr(X_{t,j} = x) = \frac{(k / R_t)^k}{x! \Gamma(k)} \int_0^\infty \nu^{(x+k)-1} \exp\left\{ -\nu \left( 1 + \frac{k}{R_t} \right) \right\} d\nu. \label{eq:integral_step2}
\end{equation}
观察被积函数指数部分，引入换元变量：
令
\begin{equation}
u \triangleq \nu \left( 1 + \frac{k}{R_t} \right) = \nu \left( \frac{k + R_t}{R_t} \right),
\end{equation}
则积分变量与微分元分别映射为：
\begin{equation}
\nu = u \left( \frac{R_t}{k + R_t} \right), \quad d\nu = \left( \frac{R_t}{k + R_t} \right) du.
\end{equation}
由于 $1 + k/R_t > 0$，当 $\nu \to 0$ 时 $u \to 0$，当 $\nu \to \infty$ 时 $u \to \infty$，积分限保持 $(0, \infty)$ 不变。将代换关系回代入式~\eqref{eq:integral_step2}：
\begin{align}
\int_0^\infty \nu^{(x+k)-1} e^{-u} d\nu &= \int_0^\infty \left[ u \left( \frac{R_t}{k + R_t} \right) \right]^{(x+k)-1} e^{-u} \cdot \left( \frac{R_t}{k + R_t} \right) du \notag \\
&= \left( \frac{R_t}{k + R_t} \right)^{x+k} \int_0^\infty u^{(x+k)-1} e^{-u} du.
\end{align}
由欧拉 Gamma 函数的严格解析定义，后项定积分严格收敛并满足：
\begin{equation}
\int_0^\infty u^{(x+k)-1} e^{-u} du \equiv \Gamma(x+k).
\end{equation}
因此，积分部分完全化简为闭式表达：
\begin{equation}
\int_0^\infty \nu^{(x+k)-1} \exp\left\{ -\nu \left( 1 + \frac{k}{R_t} \right) \right\} d\nu = \Gamma(x+k) \left( \frac{R_t}{k + R_t} \right)^{x+k}.
\end{equation}
将其与提取出的常数系数相乘：
\begin{align}
\Pr(X_{t,j} = x) &= \frac{(k / R_t)^k}{x! \Gamma(k)} \cdot \Gamma(x+k) \left( \frac{R_t}{k + R_t} \right)^{x+k} \notag \\
&= \frac{\Gamma(x+k)}{x! \Gamma(k)} \cdot \frac{k^k}{R_t^k} \cdot \frac{R_t^x R_t^k}{(k + R_t)^{x+k}} \notag \\
&= \frac{\Gamma(x+k)}{x! \Gamma(k)} \cdot \frac{k^k R_t^x}{(k + R_t)^k (k + R_t)^x}.
\end{align}
按幂次基数对分子分母进行归并：
\begin{equation}
\Pr(X_{t,j} = x) = \frac{\Gamma(x+k)}{x! \Gamma(k)} \left( \frac{k}{k + R_t} \right)^k \left( \frac{R_t}{k + R_t} \right)^x, \quad x = 0, 1, 2, \dots. \label{eq:nb_pmf_derived}
\end{equation}
式~\eqref{eq:nb_pmf_derived} 严格表明：**在个体传染力服从 Gamma 分布的微观异质性假定下，边缘化后的个体二代感染数严格服从负二项分布 $\NB(k, p)$**，其中形状参数为过度离散参数 $k$，成功概率参数定义为：
\begin{equation}
p \triangleq \frac{k}{k + R_t}. \label{eq:prob_p_def}
\end{equation}
从连续微观混合积分到离散负二项 PMF 的逐行代数证明至此完全证毕。

\subsection{成功概率 $p$ 的反直觉统计物理顿悟}
在常规的数理统计教学中，负二项分布 $\NB(r, p)$ 常被直观地定义为：“在独立重复的伯努利掷硬币试验中，累计观察到第 $r$ 次‘成功’之前，所经历的‘失败总次数’ $X$ 的概率分布”。

在流行病学分支过程的语境下，很多初学者极易产生困惑：在传染病暴发中，被感染的人数越多代表疫情越失控，为什么参数形式却反直觉地呈现为 $p = \frac{k}{k + R_t}$？究竟什么是“成功”，什么是“失败”？

理解这一参数的核心在于建立正确的**传染动力学物理顿悟**：
\begin{enumerate}[label=(\arabic*)]
    \item **“传染事件”在数学上对应一次“失败”**：
    当一个感染者接触易感者并成功将病毒传递给对方时，传播链条得到了进一步延长与扩张，这在试验模型中被记为一次试验“失败”（发生了一例二代感染 $X$）；
    \item **“传播阻断 / 链条终止”在数学上对应一次“成功”**：
    当感染者遇到具备免疫力的人群、或者因康复、隔离与死亡而未引发二次传染时，该分支链条走向完结终止，这在数学模型中被记为一次“成功”（阻断成功）；
    \item **均值方程的反解自洽性**：
    负二项分布的标准均值公式为 $\E[X] = r \frac{1-p}{p}$。在个体模型中 $r = k$，均值必须恒等于当期期望再生数 $R_t$：
    \begin{equation}
    R_t = k \frac{1-p}{p} \implies p R_t = k(1-p) = k - kp \implies p(k + R_t) = k \implies p = \frac{k}{k + R_t}.
    \end{equation}
\end{enumerate}
这一反解逻辑赋予了 $p$ 极其深刻的流行病学物理直觉：
若病毒传播力极强（$R_t \gg 1$），则终止链条的概率 $p \to 0$。这意味着传播极难被阻断，在累计达成 $k$ 次阻断（成功）之前，个体将经历大量的传染事件（失败），从而在宏观上表现为极其猛烈的超级传播暴发；
反之，若实施了严格的非药物干预（NPIs）导致 $R_t \to 0$，则阻断概率 $p \to 1$，传播链条几乎立刻终止。

\section{宏观观测系统的不完全性与二项报告抽稀过程}\label{sec:thinning}

\subsection{二项报告抽稀的形式化定义}
在真实公共卫生监测实践中，全社会底层实际发生的感染人数 $I_t$ 是无法被全量观测到的潜在潜变量。受制于无症状感染比例、轻症不就医行为以及实验室检测可及性，只有一部分感染者会前往哨点医院就诊并被国家传染病监测网络记录。

本文将这一宏观观测衰减过程形式化建模为**条件二项报告抽稀过程（Binomial Reporting Thinning Process）**：
假定在第 $t$ 代，每个实际新发感染者独立地以固定概率 $\rho \in (0, 1]$ 被医疗监测系统捕获并确诊上报。记第 $l$ 个实际感染者被上报的指示变量为 $B_l \sim \operatorname{Bernoulli}(\rho)$（$l = 1, 2, \dots, I_t$），且在给定 $I_t$ 的条件下，各个体的上报指示变量相互条件独立。

因此，疾控系统在第 $t$ 代实际记录到的报告病例数 $C_t$ 是一个复合随机和：
\begin{equation}
C_t = \sum_{l=1}^{I_t} B_l. \label{eq:binomial_thinning_sum}
\end{equation}
条件分布显然满足二项分布：
\begin{equation}
C_t \mid I_t \sim \Binomial(I_t, \rho).
\end{equation}

\subsection{报告率参数 $\rho$ 的真实流行病学折算逻辑}
在流行病学实证分析中，$\rho$ 绝非普通的实验室核酸检测灵敏度，而是**“全社会有症状感染者最终转化为实验室确诊住院病例的系统级复合报告率”**。

以本文后续实证所使用的全美 CDC NHSN 流感住院监测数据为例，参数 $\rho$ 的数量级具有坚实的临床流行病学依据：
\begin{itemize}
    \item 依据 Zhou 等人 \citep{zhou2012hospitalizations} 基于全美住院样本数据库（NIS）的长期回溯测算，美国季节性流感每年引发的实验室确诊住院人数平均处于 20 万至 40 万例区间；
    \item 依据 Tokars 等人 \citep{tokars2018seasonal} 基于前瞻性社区队列与血清抗体流调的权威测算，全美季节性流感每年的平均有症状侵袭率（Symptomatic Attack Rate）约为 8\% 至 15\%，对应全美每年约有 2,500 万至 5,000 万人发生有症状感染；
    \item 将两组权威数据相除，可严格折算出全美流感感染者转化为 NHSN 住院报告病例的复合比例区间为：
    \begin{equation}
    \rho \approx \frac{20\text{ 万}}{3500\text{ 万}} \sim \frac{40\text{ 万}}{2500\text{ 万}} \in [0.0056, \; 0.0210]. \label{eq:rho_range}
    \end{equation}
\end{itemize}
即 $\rho$ 处于千分之六至百分之二的极小数量级区间。这一现实背景凸显了二项报告抽稀模型的必要性：实际监测面临的是从浩瀚的微观感染池中进行极端稀疏的随机抽样。

\section{概率母函数复合律与分布族保族封闭性定理}\label{sec:pgf_closure}

当潜变量 $I_t$ 服从负二项分支和、而观测变量 $C_t$ 服从二项报告抽稀时，观测变量 $C_t$ 的边际分布形式是什么？是否会破坏负二项分布的数学结构？

本节利用概率母函数（Probability Generating Function, PGF）技术，严格证明其分布族的保族封闭性。

\subsection{单代独立感染者和的概率母函数}
由式~\eqref{eq:nb_pmf_derived}，单个感染者的二代感染数 $X_{t,j} \sim \NB(k, p)$，其中 $p = \frac{k}{k+R_t}$。
负二项分布的标准 PGF 解析形式为：
\begin{equation}
G_X(s) \triangleq \E[s^X] = \left( \frac{p}{1 - (1-p)s} \right)^k.
\end{equation}
根据式~\eqref{eq:branching_sum}，在给定上一代真实风险集规模 $I_{t-1} = n$ 的条件下，真实感染数 $I_t = \sum_{j=1}^n X_{t,j}$ 为 $n$ 个独立同分布随机变量之和。根据独立和的母函数性质：
\begin{equation}
G_{I_t}(s) = \prod_{j=1}^n G_X(s) = \left( G_X(s) \right)^n = \left( \frac{p}{1 - (1-p)s} \right)^{nk}. \label{eq:pgf_It}
\end{equation}
因此，当期真实感染总数满足负二项可加性：$I_t \mid n, R_t \sim \NB(nk, p)$，其形状参数累加为 $nk$。

\subsection{PGF 复合运算与分布族封闭性证明}
根据式~\eqref{eq:binomial_thinning_sum}，报告数 $C_t = \sum_{l=1}^{I_t} B_l$ 为复合随机变量和，其中 $B_l \sim \operatorname{Bernoulli}(\rho)$ 的母函数为：
\begin{equation}
G_B(s) \triangleq \E[s^B] = (1 - \rho) + \rho s.
\end{equation}
利用经典复合随机变量的母函数复合定理：
\begin{equation}
G_{C_t}(s) = G_{I_t}\big( G_B(s) \big).
\end{equation}
将式~\eqref{eq:pgf_It} 中的变量 $s$ 替换为 $G_B(s)$：
\begin{equation}
G_{C_t}(s) = \left( \frac{p}{1 - (1-p)\big[(1-\rho) + \rho s\big]} \right)^{nk}. \label{eq:pgf_Ct_step1}
\end{equation}

下面给出保姆级分母通分与公因子提取代数推导：

展开分母表达式：
\begin{align}
\text{分母} &= 1 - (1-p)(1-\rho) - \rho(1-p)s \notag \\
&= 1 - \big[ 1 - \rho - p + p\rho \big] - \rho(1-p)s \notag \\
&= \big[ p + \rho(1-p) \big] - \rho(1-p)s.
\end{align}
分子分母同时除以提取因子 $[p + \rho(1-p)]$：
\begin{equation}
G_{C_t}(s) = \left( \frac{\frac{p}{p + \rho(1-p)}}{1 - \left[ \frac{\rho(1-p)}{p + \rho(1-p)} \right] s} \right)^{nk}. \label{eq:pgf_Ct_step2}
\end{equation}
定义经过抽稀后的有效成功概率为：
\begin{equation}
p^* \triangleq \frac{p}{p + \rho(1-p)}. \label{eq:p_star_def}
\end{equation}
将原本的成功概率 $p = \frac{k}{k+R_t}$ 与失败概率 $1-p = \frac{R_t}{k+R_t}$ 代入式~\eqref{eq:p_star_def} 中进行通分化简：
\begin{align}
p^* &= \frac{\frac{k}{k + R_t}}{\frac{k}{k + R_t} + \rho \cdot \frac{R_t}{k + R_t}} \notag \\
&= \frac{\frac{k}{k + R_t}}{\frac{k + \rho R_t}{k + R_t}} \notag \\
&= \frac{k}{k + \rho R_t}. \label{eq:p_star_derived}
\end{align}
相应地，有效失败概率为：
\begin{equation}
1 - p^* = 1 - \frac{k}{k + \rho R_t} = \frac{\rho R_t}{k + \rho R_t}.
\end{equation}
将 $p^*$ 回代入式~\eqref{eq:pgf_Ct_step2}：
\begin{equation}
G_{C_t}(s) = \left( \frac{p^*}{1 - (1 - p^*) s} \right)^{nk}. \label{eq:pgf_Ct_closed}
\end{equation}
对比标准负二项分布母函数形式，式~\eqref{eq:pgf_Ct_closed} 在代数上严格证明了以下定理：
\begin{theorem}[\textbf{二项报告抽稀保族封闭性定理}]\label{thm:closure}
设风险集规模为 $n$，个体传染力过度离散参数为 $k$，当期瞬时再生数为 $R_t$。若微观传播服从负二项分支过程，宏观报告服从条件独立二项报告抽稀过程，则观测病例数 $C_t$ 严格服从形状参数为 $nk$、成功概率为 $p^* = \frac{k}{k+\rho R_t}$ 的负二项分布族闭包：
\begin{equation}
C_t \mid n, R_t \sim \NB\left( nk, \; \frac{k}{k + \rho R_t} \right).
\end{equation}
\end{theorem}

该定理具有极高的理论优雅性：二项报告抽稀并不会破坏负二项分布的族系属性，而是将有效均值参数严格收缩为 $\rho R_t$，同时完好保留了形状参数 $nk$。

\subsection{观测变量的二阶矩结构解析}
由负二项分布的标准矩性质，观测病例数 $C_t$ 的条件期望与条件方差可显式计算如下：
\begin{equation}
\E[C_t \mid n, R_t] = nk \cdot \frac{1 - p^*}{p^*} = nk \cdot \frac{\frac{\rho R_t}{k + \rho R_t}}{\frac{k}{k + \rho R_t}} = nk \cdot \frac{\rho R_t}{k} = \rho n R_t.
\end{equation}
定义上一代风险集 $n$ 对应的期望被报告人数为 $m \triangleq \rho n$，则：
\begin{equation}
\E[C_t \mid n, R_t] = m R_t.
\end{equation}

进一步展开条件方差：
\begin{equation}
\Var(C_t \mid n, R_t) = nk \cdot \frac{1 - p^*}{(p^*)^2} = nk \cdot \frac{\frac{\rho R_t}{k + \rho R_t}}{\left( \frac{k}{k + \rho R_t} \right)^2} = nk \cdot \frac{\rho R_t}{k + \rho R_t} \cdot \frac{(k + \rho R_t)^2}{k^2}.
\end{equation}
分子分母对消公因子 $k(k + \rho R_t)$：
\begin{align}
\Var(C_t \mid n, R_t) &= \frac{n \rho R_t (k + \rho R_t)}{k} \notag \\
&= \frac{n \rho k R_t + n \rho^2 R_t^2}{k} \notag \\
&= \rho n R_t + \frac{(\rho R_t)^2 n}{k}. \label{eq:obs_var_decomp}
\end{align}
式~\eqref{eq:obs_var_decomp} 极其清晰地展现了二阶矩结构的数理正交叠加：
\begin{itemize}
    \item 第一项 $\rho n R_t = m R_t$：纯泊松随机相遇计数波动（Poisson Counting Noise）；
    \item 第二项 $\frac{(\rho R_t)^2 n}{k} = \frac{\rho m R_t^2}{k}$：个体生物学生殖数过度离散引发的超传播乘性放大项（Superspreading Inflation Noise）。
\end{itemize}

\section{微观分支过程与宏观确定性仓室 ODE 的流体极限动力学桥梁}\label{sec:fluid_limit}

为了深化论文的动力系统机理基础，并架起微观离散统计推断与马知恩先生学派经典常微分方程动力系统的内在血脉连接，本节系统阐明微观随机分支过程与宏观仓室 ODE 之间的 Kurtz 流体极限对应关系。

\subsection{经典 Kermack--McKendrick SIR 动力系统与相空间稳定性}
在经典数理流行病学中，Kermack 与 McKendrick \citep{kermack1927contribution} 以及马知恩等 \citep{ma2001stability,ma2004epidemic} 建立的确定性无时滞 SIR 仓室微分方程系统具有如下标准形式：
\begin{equation}
\begin{cases}
\dfrac{dS}{dt} = -\dfrac{\beta S I}{N}, \\[2ex]
\dfrac{dI}{dt} = \dfrac{\beta S I}{N} - \gamma I, \\[2ex]
\dfrac{dR}{dt} = \gamma I,
\end{cases} \label{eq:sir_ode}
\end{equation}
其中 $N = S(t) + I(t) + R(t)$ 为封闭总人口规模，$\beta$ 为有效接触传播率，$\gamma$ 为传染期恢复率。基本再生数定义为无量纲参数 $\mathcal{R}_0 \triangleq \beta / \gamma$。

在常微分方程定性理论中，系统~\eqref{eq:sir_ode} 恒存在无病平衡点 $E_0 = (N, 0, 0)$。对其进行线性化分析，计算系统 Jacobian 矩阵在 $E_0$ 处的特征值：
\begin{equation}
J(E_0) = \begin{pmatrix}
0 & -\beta & 0 \\
0 & \beta - \gamma & 0 \\
0 & \gamma & 0
\end{pmatrix}.
\end{equation}
显然，决定系统稳定性的关键特征根为 $\lambda = \beta - \gamma = \gamma(\mathcal{R}_0 - 1)$：
\begin{itemize}
    \item 当 $\mathcal{R}_0 < 1$ 时，$\lambda < 0$，无病平衡点 $E_0$ 局部渐近稳定，疫情自发衰退；
    \item 当 $\mathcal{R}_0 > 1$ 时，$\lambda > 0$，$E_0$ 失稳，系统在相空间中发生经典的**跨临界分岔（Transcritical Bifurcation）**，疫情呈指数暴发。
\end{itemize}

\subsection{纯跳跃马尔可夫生灭过程的生成元与鞅表示}
在有限微观人群中，疾病传播本质是一个连续时间纯跳跃马尔可夫链。设在时刻 $t$，系统状态为 $(S_N(t), I_N(t))$。系统可能的状态转移及其跳跃速率严格定义为：
\begin{align}
(S, I) \to (S-1, I+1) \quad &\text{转移速率：} \quad q_1(S, I) = \frac{\beta S I}{N}, \\
(S, I) \to (S, I-1) \quad &\text{转移速率：} \quad q_2(S, I) = \gamma I.
\end{align}
根据马尔可夫过程的半鞅分解定理（Doob--Meyer 分解），连续时间跳跃过程可严格表示为漂移测度积分与零均值局部平方可积鞅之和：
\begin{equation}
\frac{I_N(t)}{N} = \frac{I_N(0)}{N} + \int_0^t \left( \beta \frac{S_N(u)}{N} \frac{I_N(u)}{N} - \gamma \frac{I_N(u)}{N} \right) du + M_N(t), \label{eq:semimartingale}
\end{equation}
其中 $M_N(t)$ 为二次可积鞅，其尖括号可预测变差过程（Predictable Quadratic Variation）满足：
\begin{equation}
\langle M_N \rangle_t = \frac{1}{N} \int_0^t \left( \beta \frac{S_N(u)}{N} \frac{I_N(u)}{N} + \gamma \frac{I_N(u)}{N} \right) du. \label{eq:martingale_variation}
\end{equation}

\subsection{Kurtz 密度相关大数定律（Fluid Limit Theorem）}
马氏过程与微分方程定性理论的终极汇聚点在于 Thomas Kurtz 于 1970 年证明的著名极限定理 \citep{kurtz1970solutions,kurtz1971limit,ethier1986markov}：
\begin{theorem}[\textbf{Kurtz 密度相关大数定律}]\label{thm:kurtz}
设宏观总人口规模 $N \to \infty$。若初始归一化状态几乎必然收敛：
\begin{equation}
\lim_{N \to \infty} \left( \frac{S_N(0)}{N}, \frac{I_N(0)}{N} \right) = (s_0, i_0),
\end{equation}
则在任意有限时间区间 $[0, T]$ 内，微观纯跳跃马尔可夫随机过程的密度路径依概率 1 一致收敛于确定性常微分方程系统~\eqref{eq:sir_ode} 的解轨道 $(s(t), i(t))$：
\begin{equation}
\sup_{0 \le t \le T} \left\| \left( \frac{S_N(t)}{N}, \frac{I_N(t)}{N} \right) - (s(t), i(t)) \right\| \xrightarrow[N \to \infty]{\text{a.s.}} 0.
\end{equation}
\end{theorem}
证明的核心由式~\eqref{eq:martingale_variation} 直接导出：因为鞅变差包含 $1/N$ 前置系数，当 $N \to \infty$ 时，随机涨落项 $\langle M_N \rangle_t \to 0$，鞅项以概率 1 溃缩为零，随机动力系统完全退化为确定性均值向量场流形。

\subsection{动力学理论定位与本文研究切入点}
Kurtz 流体极限理论为本文的研究切入点提供了极其崇高的数理定位：
\begin{center}
\fbox{\parbox{0.92\textwidth}{
\textbf{确定性常微分方程仓室系统，本质上是大数极限 $N \to \infty$ 下微观随机过程的均值流形；\\
而真实流行病学监测的核心挑战恰恰在于：现实中的监测样本规模并非无穷大（$n < \infty, m < \infty$）。本文所聚焦研究的，正是微观随机涨落偏离确定性动力学轨道的方差标度律与最优抽样边界！}
}}
\end{center}
这一动力学桥梁的建立，彻底打通了微观分支随机推断与宏观微分方程定性理论的学科阻隔。在下一章中，我们将正式推导以不同推断目标为导向的渐近方差分解与运筹设计理论。


\section{呼吸道传染病过度离散参数 $k$ 经验校准综述}\label{sec:k_calibration}

在第二章推导的负二项分支过程模型中，生物学生殖数过度离散参数 $k$ 扮演着决定系统微观方差乘性放大强度的关键角色。为了让本文的理论分析具备扎实的生物学与流行病学先验基础，本节对近二十年来国际顶级期刊中关于典型人类传染病离散参数 $k$ 的实证测量文献进行系统的元分析梳理。

\subsection{主要病原体离散参数测量大表}
表~\ref{tab:superspreading_meta} 汇总了涵盖冠状病毒属、正粘病毒属、副粘病毒属等全球主要呼吸道及接触传播病原体的经验测定参数。

\begin{table}[htbp]
\centering
\small
\caption{全球主要人类传染病生殖数过度离散参数 $k$ 实证测量元分析汇总表}
\label{tab:superspreading_meta}
\begin{tabularx}{\textwidth}{lccccX}
\toprule
\textbf{病原体与流行事件} & \textbf{典型 $R_0$ 区间} & \textbf{离散参数 $k$} & \textbf{20/80 集中度} & \textbf{超传播特征定性} & \textbf{权威文献来源} \\
\midrule
\textbf{SARS-CoV-1} (2003) & $2.2 \sim 3.6$ & $\mathbf{0.16}$ & 80\% 来自 9\% 个体 & 极端强超传播 & Lloyd-Smith et al. (2005) \citep{lloydsmith2005superspreading} \\
\textbf{SARS-CoV-2} (原始株) & $2.5 \sim 3.0$ & $\mathbf{0.10 \sim 0.30}$ & 80\% 来自 10\% 个体 & 强超传播 & Endo et al. (2020) \citep{endo2020estimating} \\
\textbf{SARS-CoV-2} (Delta/Omicron) & $5.0 \sim 9.5$ & $\mathbf{0.25 \sim 0.45}$ & 80\% 来自 15\% 个体 & 中强超传播 & Althouse et al. (2020) \citep{althouse2020superspreading} \\
\textbf{MERS-CoV} (2012) & $0.4 \sim 0.9$ & $\mathbf{0.26}$ & 局部医院暴发 & 强超传播 & Kucharski et al. (2015) \\
\textbf{季节性流感} (Influenza A/H3N2) & $1.2 \sim 1.6$ & $\mathbf{0.35 \sim 0.55}$ & 80\% 来自 22\% 个体 & 中度过度离散 & Fraser et al. (2007) \citep{fraser2007estimating} \\
\textbf{大流行流感} (2009 H1N1pdm) & $1.4 \sim 1.8$ & $\mathbf{0.32 \sim 0.48}$ & 80\% 来自 20\% 个体 & 中度过度离散 & White et al. (2009) \citep{white2009estimation} \\
\textbf{麻疹} (Measles, 历史株) & $12.0 \sim 18.0$ & $\mathbf{0.10 \sim 0.15}$ & 高度集中暴发 & 强超传播 & Anderson \& May (1991) \citep{anderson1991infectious} \\
\textbf{天花} (Smallpox, 历史株) & $3.5 \sim 6.0$ & $\mathbf{0.32 \sim 0.60}$ & 局部密切接触 & 中度过度离散 & Gani \& Leach (2001) \\
\textbf{肺炎双球菌} (Pneumococcal) & $1.1 \sim 1.4$ & $\mathbf{1.20 \sim 2.50}$ & 接近泊松分布 & 弱过度离散 & Lloyd-Smith et al. (2005) \citep{lloydsmith2005superspreading} \\
\bottomrule
\end{tabularx}
\end{table}

\subsection{对本文参数设定的启示}
从表~\ref{tab:superspreading_meta} 可以清晰看出两个显著特征：
\begin{enumerate}[label=(\arabic*)]
    \item 呼吸道病原体普遍处于 $k \in (0.1, 0.6)$ 的区间之内，经典泊松分布假定（$k = \infty$）在现实中根本不存在；
    \item 本文在第四章数值模拟中选取 $k = 0.35$ 作为流感基准，并在第五章唯象实证中覆盖 $k \in [0.2, 0.5]$，完全处于季节性流感与新型流感的临床测量中心区间，具有坚实的事实依据。
\end{enumerate}
